\section{Conclusiones}


A partir de las experiencias realizadas en el laboratorio con el sensor de temperatura NTC, se verificó de manera experimental la relación inversamente proporcional y no lineal entre la resistencia eléctrica del dispositivo semiconductor y la temperatura a la que se encuentra expuesto, tal como postula el modelo teórico de Steinhart-Hart. Asimismo, las pruebas desarrolladas permitieron identificar la notable sensibilidad del componente ante variaciones térmicas bruscas, lo cual demanda métodos de medición rigurosos y estables.

En relación con los métodos de adquisición empleados, la utilización del sistema microcontrolado mediante la placa Arduino UNO y la aplicación de los coeficientes provistos por el fabricante demostró ser la metodología más confiable y precisa para la determinación de la temperatura. Aunque las mediciones manuales con multímetro permitieron observar la tendencia decreciente general, sufrieron un alto grado de incertidumbre debido a la fluctuación constante en las lecturas, la falta de acoplamiento mecánico firme en las terminales y la transferencia de calor indeseada durante la manipulación de las sondas, afectando sustancialmente la trazabilidad de los datos.

Por otra parte, se evidenció la crítica importancia de contar con coeficientes de calibración precisos. El intento de recalcular los parámetros de Steinhart-Hart a partir de mediciones manuales arrojó resultados altamente inconsistentes y sin sentido físico, lo que pone de manifiesto la extrema sensibilidad de los términos en el denominador de dicha ecuación. La propagación de pequeños errores en las lecturas primarias de temperatura y resistencia, sumada a la aproximación en la conversión de unidades, distorsiona gravemente el modelo matemático del sensor.

En conclusión, para garantizar la validez en aplicaciones de instrumentación biomédica donde se emplean termistores NTC, no solo es indispensable la correcta calibración e inclusión de los coeficientes originales del fabricante, sino también el diseño de un acondicionamiento de señal adecuado. Esto implica implementar conexiones mecánicas y eléctricas estables, minimizar el ruido mediante filtrado adecuado y contemplar los retardos dinámicos del conversor analógico-digital para obtener mediciones reproducibles y con alta fidelidad clínica.

